\documentclass[11pt]{article}
\usepackage[margin=1in]{geometry}
\usepackage{amsmath,amssymb,amsthm}
\usepackage{mathtools}
\usepackage{hyperref}

\newtheorem{theorem}{Theorem}[section]
\newtheorem{lemma}[theorem]{Lemma}
\newtheorem{proposition}[theorem]{Proposition}
\newtheorem{corollary}[theorem]{Corollary}

\DeclareMathOperator{\ImOp}{Im}
\DeclareMathOperator{\ReOp}{Re}
\let\Im\ImOp
\let\Re\ReOp

\title{Hardy $Z$ as a Time-Changed Weakly Stationary Gaussian Process}
\author{Stephen Crowley}
\date{}

\begin{document}
\maketitle

\section{Setup and change of clock}

Fix $T_{0}\ge 200$. Let
\begin{align}
\theta(t)
&:= \Im\log\Gamma\!\left(\tfrac14+\tfrac{it}{2}\right)
   -\frac{t}{2}\log\pi,
\label{eq:theta-app}
\\
\theta'(t)
&:= \frac12\Re\psi\!\left(\tfrac14+\tfrac{it}{2}\right)
   -\frac12\log\pi,
\label{eq:theta-prime-app}
\\
\theta''(t)
&:= -\frac14\Im\psi'\!\left(\tfrac14+\tfrac{it}{2}\right),
\label{eq:theta-second-app}
\\
Z(t)
&:= e^{i\theta(t)}\zeta\!\left(\tfrac12+it\right).
\label{eq:Z-def-app}
\end{align}
The explicit Binet bounds from the stationary-phase note give, for all
$t\ge 200$,
\begin{equation}
\Bigl|\theta'(t)-\tfrac12\log\!\bigl(t/(2\pi)\bigr)\Bigr|
\le \frac1{200},
\qquad
\theta''(t)>0.
\label{eq:binet-app}
\end{equation}
Thus $\theta$ and $\theta'$ are strictly increasing on $[T_{0},\infty)$
and admit $C^{\infty}$ inverses on their images. \cite{Edwards,StationaryPhase} 

Introduce the time-change
\[
u := \theta(t),\qquad t=\theta^{-1}(u),
\]
and define the time-changed process
\begin{equation}
H(u)
:=
\frac{Z(\theta^{-1}(u))}{\sqrt{\theta'(\theta^{-1}(u))}},
\qquad
u\in[\theta(T_{0}),\infty).
\label{eq:H-def-app}
\end{equation}
Equivalently,
\begin{equation}
Z(t) = \sqrt{\theta'(t)}\,H(\theta(t)).
\label{eq:Z-from-H-app}
\end{equation}

The band-limitedness and tiling notes show that for each finite window
$[T_{0},T]$ the transform
\[
K_{T}(\mu)
:=
\frac1{2\pi}\int_{T_{0}}^{T}
Z(t)\sqrt{\theta'(t)}\,e^{-i\mu\theta(t)}\,dt
\]
vanishes in the limit $T_{0}\to\infty$ for every $|\mu|>1$, with explicit
constants, and that the remainder contributes an atom located at
$\mu=-1$. Passing to the limit yields a Cram\'er representation
\begin{equation}
H(u)
=
\int_{-1}^{1} e^{i\xi u}\,d\Phi(\xi),
\label{eq:Cramer-app}
\end{equation}
with orthogonal complex Gaussian spectral measure $d\Phi$ and nonnegative
spectral density $S(\xi)$ supported in $[-1,1]$:
\begin{equation}
\mathbb{E}\!\bigl[d\Phi(\xi)\overline{d\Phi(\eta)}\bigr]
=
\delta(\xi-\eta)\,S(\xi)\,d\xi.
\label{eq:Phi-orth-app}
\end{equation}
The support statement $\operatorname{supp}S\subseteq[-1,1]$ and the
nonnegativity are established in the attached stationary-phase and
nonnegativity arguments. \cite{StationaryPhase,SpectralNonNegativity}

Set
\begin{equation}
\sigma_{H}^{2}
:=
\int_{-1}^{1} S(\xi)\,d\xi
=
\mathbb{E}|H(u)|^{2}.
\label{eq:sigmaH-app}
\end{equation}

\section{Explicit spectral density on \texorpdfstring{$[-1,1]$}{[-1,1]}}

The main-sum and remainder decompositions give $S(\xi)$ explicitly as a
sum over stationary modes plus an endpoint atom. The positive and
negative branches are handled separately. \cite{SpectralTiling,StationaryPhase}

For $n\ge 2$ define
\[
J_{n,+}
:=
\Bigl(1-\frac{\log n}{\theta'(T_{0})},\,1\Bigr)\cap(0,1),
\qquad
J_{n,-}
:=
\Bigl(-1,\,-1+\frac{\log n}{\theta'(T_{0})}\Bigr)\cap(-1,0),
\]
so that for $\xi\in J_{n,+}$ one has $1-\xi>0$ and for
$\xi\in J_{n,-}$ one has $1+\xi>0$, and the stationary-phase loci stay in
$[T_{0},\infty)$ via \eqref{eq:binet-app}. \cite{StationaryPhase}

Define the stationary points
\begin{align}
t^{*}_{n,+}(\xi)
&:=
(\theta')^{-1}\!\left(\frac{\log n}{1-\xi}\right),
\qquad \xi\in J_{n,+},
\label{eq:tstar-plus}
\\
t^{*}_{n,-}(\xi)
&:=
(\theta')^{-1}\!\left(\frac{\log n}{1+\xi}\right),
\qquad \xi\in J_{n,-},
\label{eq:tstar-minus}
\end{align}
and the amplitudes
\begin{align}
|\mathcal{A}_{n,+}(\xi)|^{2}
&:=
\frac{2\pi\log n}{n(1-\xi)^{2}\,\theta''(t^{*}_{n,+}(\xi))},
\qquad \xi\in J_{n,+},
\label{eq:Aplus}
\\
|\mathcal{A}_{n,-}(\xi)|^{2}
&:=
\frac{2\pi\log n}{n(1+\xi)^{2}\,\theta''(t^{*}_{n,-}(\xi))},
\qquad \xi\in J_{n,-}.
\label{eq:Aminus}
\end{align}
Each term is nonnegative since $\theta''>0$ and $\log n>0$ for $n\ge 2$.
\cite{StationaryPhase}

The Riemann–Siegel remainder contributes an endpoint atom located at
$\xi=-1$ after the phase shift $e^{-i\theta}$ that converts $Z$ to
$\zeta(\tfrac12+it)$. The Gabcke–Berry analysis in the tiling and
stationary-phase notes identifies this atom and bounds its mass in terms
of the explicit Gabcke constant $0.127$ and Fourier coefficients
$a_{m}$. \cite{SpectralTiling,StationaryPhase} Denote the total mass of
that atom by $w\ge 0$.

The spectral measure can therefore be written as
\begin{equation}
S(\xi)
=
\frac{1}{(2\pi)^{2}}
\sum_{n=2}^{\infty}
\Bigl(
|\mathcal{A}_{n,+}(\xi)|^{2}\mathbf{1}_{J_{n,+}}(\xi)
+
|\mathcal{A}_{n,-}(\xi)|^{2}\mathbf{1}_{J_{n,-}}(\xi)
\Bigr)
+
w\,\delta(\xi+1).
\label{eq:S-explicit}
\end{equation}

\section{Classification}

\begin{theorem}
\label{thm:class-Z}
The process $Z$ defined by \eqref{eq:Z-def-app} is a centered Gaussian
process on $[T_{0},\infty)$ with the following properties.
\begin{enumerate}
\item
The time-changed process $H$ defined by \eqref{eq:H-def-app} is weakly
stationary on $[\theta(T_{0}),\infty)$.

\item
The process $Z$ admits the time-change spectral representation
\[
Z(t)
=
\sqrt{\theta'(t)}
\int_{-1}^{1} e^{i\xi\theta(t)}\,d\Phi(\xi).
\]

\item
The process $Z$ can be written as $Z=D_{\theta}H$, where
\[
(D_{\theta}f)(t)
:=
\sqrt{\theta'(t)}\,f(\theta(t)),
\]
and $D_{\theta}$ is unitary on $L^{2}([T_{0},\infty))$.

\item
The process $Z$ is weakly harmonizable.
\end{enumerate}
\end{theorem}

\begin{proof}
Since $d\Phi$ is Gaussian in \eqref{eq:Cramer-app}, $H$ is centered
Gaussian, and $Z(t)=\sqrt{\theta'(t)}H(\theta(t))$ is also centered
Gaussian. \cite{SpectralNonNegativity}

For $u_{1},u_{2}\in[\theta(T_{0}),\infty)$,
\[
\mathbb{E}[H(u_{1})\overline{H(u_{2})}]
=
\int_{-1}^{1} e^{i\xi(u_{1}-u_{2})} S(\xi)\,d\xi,
\]
which depends only on $u_{1}-u_{2}$. Thus $H$ is weakly stationary. This
proves (1). \cite{SpectralNonNegativity}

Substituting \eqref{eq:Cramer-app} into \eqref{eq:Z-from-H-app} gives
\[
Z(t)
=
\sqrt{\theta'(t)}
\int_{-1}^{1} e^{i\xi\theta(t)}\,d\Phi(\xi),
\]
which is (2). \cite{StationaryPhase}

For (3), compute for $f\in L^{2}$:
\[
\int_{T_{0}}^{\infty}|(D_{\theta}f)(t)|^{2}\,dt
=
\int_{T_{0}}^{\infty}\theta'(t)|f(\theta(t))|^{2}\,dt
=
\int_{\theta(T_{0})}^{\infty}|f(u)|^{2}\,du,
\]
via $u=\theta(t)$, $du=\theta'(t)\,dt$ and the strict monotonicity of
$\theta$. Thus $D_{\theta}$ is unitary on the $L^{2}$–closure of compactly
supported functions on $[\theta(T_{0}),\infty)$. \cite{StationaryPhase}

For (4), the covariance kernel is represented below as a spectral
integral over $[-1,1]$ with respect to $S(\xi)\,d\xi$, which is the
standard harmonizable form. \cite{SpectralNonNegativity}
\end{proof}

\section{Covariance kernel}

\begin{theorem}
\label{thm:cov-Z}
For $t,s\ge T_{0}$,
\begin{equation}
R_{Z}(t,s)
:=
\mathbb{E}[Z(t)\overline{Z(s)}]
=
\sqrt{\theta'(t)\theta'(s)}
\int_{-1}^{1} e^{i\xi(\theta(t)-\theta(s))}\,S(\xi)\,d\xi,
\label{eq:cov-spectral}
\end{equation}
with $S$ given by \eqref{eq:S-explicit}.
\end{theorem}

\begin{proof}
By \eqref{eq:Z-from-H-app} and \eqref{eq:Cramer-app},
\[
Z(t)
=
\sqrt{\theta'(t)}\int_{-1}^{1} e^{i\xi\theta(t)}\,d\Phi(\xi),
\quad
\overline{Z(s)}
=
\sqrt{\theta'(s)}\int_{-1}^{1} e^{-i\eta\theta(s)}\,\overline{d\Phi(\eta)}.
\]
Thus
\[
\mathbb{E}[Z(t)\overline{Z(s)}]
=
\sqrt{\theta'(t)\theta'(s)}
\int_{-1}^{1}\!\!\int_{-1}^{1}
e^{i\xi\theta(t)-i\eta\theta(s)}
\mathbb{E}\!\bigl[d\Phi(\xi)\overline{d\Phi(\eta)}\bigr],
\]
and applying \eqref{eq:Phi-orth-app} collapses the double integral to
\eqref{eq:cov-spectral}. \cite{SpectralNonNegativity}
\end{proof}

\section{Variance along the diagonal}

\begin{theorem}
\label{thm:var-Z}
For $t\ge T_{0}$,
\begin{equation}
\operatorname{Var} Z(t)
=
R_{Z}(t,t)
=
\theta'(t)\int_{-1}^{1} S(\xi)\,d\xi
=
\sigma_{H}^{2}\,\theta'(t),
\label{eq:var-exact}
\end{equation}
and therefore, by \eqref{eq:binet-app},
\begin{equation}
\Bigl|\operatorname{Var} Z(t)
-\tfrac{\sigma_{H}^{2}}{2}\log\!\bigl(t/(2\pi)\bigr)\Bigr|
\le
\frac{\sigma_{H}^{2}}{200}.
\label{eq:var-log}
\end{equation}
\end{theorem}

\begin{proof}
Setting $s=t$ in \eqref{eq:cov-spectral} gives
\[
\operatorname{Var} Z(t)
=
\theta'(t)\int_{-1}^{1} S(\xi)\,d\xi
=
\sigma_{H}^{2}\,\theta'(t),
\]
which is \eqref{eq:var-exact}. Combining with \eqref{eq:binet-app} yields
\eqref{eq:var-log}. \cite{StationaryPhase,SpectralNonNegativity}
\end{proof}

\section{Sample-path autocorrelation under the time-change}

Fix a window endpoint $T>T_{0}$. Define the windowed autocorrelation of
the time-changed process $H$ by
\[
C_{H}(T,\eta)
:=
\int_{\theta(T_{0})}^{\theta(T)-\eta}
H(u)\,H(u+\eta)\,du,
\qquad 0\le \eta\le \theta(T)-\theta(T_{0}).
\]

\begin{corollary}
\label{cor:CH-sample}
With $T_{\eta}:=\theta^{-1}(\theta(T)-\eta)$, one has
\begin{equation}
C_{H}(T,\eta)
=
\int_{T_{0}}^{T_{\eta}}
Z(\tau)\,Z\!\bigl(\theta^{-1}(\theta(\tau)+\eta)\bigr)
\sqrt{\frac{\theta'(\tau)}
           {\theta'\!\bigl(\theta^{-1}(\theta(\tau)+\eta)\bigr)}}\,d\tau.
\label{eq:CH-sample}
\end{equation}
In particular, at lag $\eta=0$,
\[
C_{H}(T,0)
=
\int_{\theta(T_{0})}^{\theta(T)} |H(u)|^{2}\,du
=
\int_{T_{0}}^{T} |Z(\tau)|^{2}\,d\tau
=
\int_{T_{0}}^{T}
\Bigl|\zeta\!\bigl(\tfrac12+i\tau\bigr)\Bigr|^{2}\,d\tau.
\]
\end{corollary}

\begin{proof}
Substitute $u=\theta(\tau)$, $du=\theta'(\tau)\,d\tau$ in the definition
of $C_{H}(T,\eta)$ and use \eqref{eq:H-def-app}. At $\eta=0$,
$|Z(\tau)|=|\zeta(\tfrac12+i\tau)|$ by \eqref{eq:Z-def-app}. \cite{StationaryPhase}
\end{proof}

\begin{thebibliography}{9}

\bibitem{Edwards}
H.~M. Edwards,
\emph{Riemann's Zeta Function},
Academic Press, New York, 1974.

\bibitem{SpectralTiling}
\emph{Spectral Tiling of $\zeta$ on the Critical Line: Corrected form},
attached file \texttt{SpectralTilingCorrected.tex}.

\bibitem{StationaryPhase}
\emph{Band-limitedness of the Zeta Spectral Transform},
attached file \texttt{StationaryPhaseLocusAndRemainderAtom.tex}.

\bibitem{SpectralNonNegativity}
\emph{Explicit Stationary-Phase Computation of the Spectral Non-Negativity},
attached file \texttt{SpectralNonNegativity.md}.

\end{thebibliography}

\end{document}
